\documentclass[12pt,a4paper]{article}
\usepackage[utf8]{inputenc}
\usepackage[T1]{fontenc}
\usepackage[french]{babel}
\usepackage{amsmath,amssymb}
\usepackage{geometry}
\geometry{margin=1.8cm}
\usepackage{enumitem}
\usepackage{parskip}
\usepackage{array}
\usepackage{booktabs}
\usepackage{eurosym}
\usepackage{xcolor}
\usepackage{tcolorbox}
\tcbuselibrary{skins,breakable}
\usepackage{tikz}
\usepackage{multicol}

% --- Couleurs du jeu ---
\definecolor{bleufonce}{RGB}{20,60,120}
\definecolor{rouge}{RGB}{180,30,30}
\definecolor{vert}{RGB}{20,130,70}
\definecolor{orange}{RGB}{230,120,20}
\definecolor{violet}{RGB}{110,30,140}
\definecolor{grisclair}{RGB}{240,240,240}

% --- Boîtes stylisées ---
\newtcolorbox{regle}[1][]{
  colback=grisclair, colframe=bleufonce, boxrule=1pt,
  arc=3mm, left=3mm, right=3mm, top=2mm, bottom=2mm, #1
}

% carte : #1 = couleur du cadre, #2 = titre
\newtcolorbox{carte}[2]{
  colback=white, colframe=#1, boxrule=1.2pt,
  arc=2mm, left=3mm, right=3mm, top=2mm, bottom=2mm,
  title=#2, fonttitle=\bfseries
}

\newtcolorbox{reponse}[1][]{
  colback=vert!5, colframe=vert, boxrule=1pt,
  arc=2mm, left=3mm, right=3mm, top=2mm, bottom=2mm, #1
}

\begin{document}

\begin{center}
{\Huge\bfseries\color{bleufonce} Le Grand Défi des Évolutions}\\[0.3cm]
{\large Jeu mathématique -- Durée : 55 minutes}\\[0.2cm]
{\large Niveau : Première / Terminale -- Taux, coefficients multiplicateurs, taux global et moyen}
\end{center}

\vspace{0.3cm}

% ============================================================
\section*{Règles du jeu}

\begin{regle}
\textbf{But du jeu :} accumuler un maximum de points en résolvant des défis sur les évolutions.\\[0.2cm]
\textbf{Organisation :} la classe est divisée en \textbf{4 équipes}. Chaque équipe désigne un capitaine qui rend les réponses.\\[0.2cm]
\textbf{Matériel :} une fiche de score par équipe, des cartes-questions, une calculatrice autorisée.\\[0.2cm]
\textbf{Déroulement :} 5 manches chronométrées. À chaque manche, les équipes marquent des points selon le barème indiqué.\\[0.2cm]
\textbf{Victoire :} l’équipe qui totalise le plus de points à la fin des 5 manches remporte le défi.
\end{regle}

\vspace{0.3cm}

\begin{center}
\renewcommand{\arraystretch}{1.4}
\begin{tabular}{|c|l|c|c|}
\hline
\textbf{Manche} & \textbf{Titre} & \textbf{Durée} & \textbf{Points max} \\
\hline
1 & Échauffement éclair & 10 min & 8 pts \\
\hline
2 & Défis en équipe & 15 min & 12 pts \\
\hline
3 & Duel de capitaines & 10 min & 8 pts \\
\hline
4 & Enveloppe mystère & 10 min & 12 pts \\
\hline
5 & Finale « Tout ou rien » & 10 min & 10 pts \\
\hline
\textbf{Total} & & \textbf{55 min} & \textbf{50 pts} \\
\hline
\end{tabular}
\end{center}

\newpage

% ============================================================
\section*{Manche 1 -- Échauffement éclair (10 min, 8 pts)}

\textit{Chaque bonne réponse rapporte 1 point. Les 8 questions sont posées à toute la classe. La première équipe qui lève la main et donne la bonne réponse marque le point.}

\vspace{0.3cm}

\begin{carte}{bleufonce}{Question 1}
Un prix augmente de \(25\%\). Quel est le coefficient multiplicateur ?
\end{carte}

\begin{carte}{bleufonce}{Question 2}
Un prix baisse de \(10\%\). Quel est le coefficient multiplicateur ?
\end{carte}

\begin{carte}{bleufonce}{Question 3}
Le coefficient multiplicateur est \(1{,}4\). Quel est le taux d’évolution en \% ?
\end{carte}

\begin{carte}{bleufonce}{Question 4}
Le coefficient multiplicateur est \(0{,}75\). Quel est le taux d’évolution en \% ?
\end{carte}

\begin{carte}{bleufonce}{Question 5}
Un article passe de \(50\)~\euro{} à \(60\)~\euro{}. Quel est le taux d’évolution ?
\end{carte}

\begin{carte}{bleufonce}{Question 6}
Une grandeur augmente de \(20\%\), puis de \(20\%\). Le taux global est-il \(40\%\) ?
\end{carte}

\begin{carte}{bleufonce}{Question 7}
Que vaut \(1{,}10 \times 0{,}90\) ?
\end{carte}

\begin{carte}{bleufonce}{Question 8}
Un prix baisse de \(50\%\). Par quel coefficient faut-il multiplier pour revenir au prix initial ?
\end{carte}

\vspace{0.3cm}
\textbf{Réponses manche 1 :} 1) \(1{,}25\) — 2) \(0{,}90\) — 3) \(+40\%\) — 4) \(-25\%\) — 5) \(+20\%\) — 6) Non, \(1{,}44\) soit \(+44\%\) — 7) \(0{,}99\) — 8) Par \(2\) (car \(0{,}5 \times 2 = 1\)).

\newpage

% ============================================================
\section*{Manche 2 -- Défis en équipe (15 min, 12 pts)}

\textit{Chaque équipe reçoit les 3 défis ci-dessous. Elle dispose de 15 minutes pour les résoudre. Chaque défi est noté sur 4 points (2 pts pour la démarche, 2 pts pour le résultat).}

\vspace{0.3cm}

\begin{carte}{rouge}{Défi A}
Un smartphone coûte \(600\)~\euro{}. Son prix baisse de \(15\%\), puis augmente de \(10\%\).\\[0.2cm]
\textbf{Quel est son prix final ? Le prix a-t-il globalement augmenté ou baissé ?}
\end{carte}

\begin{carte}{rouge}{Défi B}
Une ville comptait \(25\,000\) habitants en 2015. Elle a perdu \(4\%\) d’habitants en 2016, puis gagné \(6\%\) en 2017, puis perdu \(2\%\) en 2018.\\[0.2cm]
\textbf{Combien d’habitants compte-t-elle fin 2018 (arrondir à l’unité) ? Quel est le taux global d’évolution ?}
\end{carte}

\begin{carte}{rouge}{Défi C}
Un article subit deux évolutions successives : \(+30\%\) puis \(-30\%\).\\[0.2cm]
\textbf{Le prix final est-il égal au prix initial ? Justifier par un calcul.}
\end{carte}

\vspace{0.3cm}
\textbf{Corrigé manche 2 :}
\begin{itemize}
    \item \textbf{Défi A :} \(600 \times 0{,}85 \times 1{,}10 = 561\)~\euro{}. Baisse globale (CM \(= 0{,}935\), soit \(-6{,}5\%\)).
    \item \textbf{Défi B :} \(25\,000 \times 0{,}96 \times 1{,}06 \times 0{,}98 \approx 24\,953\) habitants. CM \(= 0{,}96 \times 1{,}06 \times 0{,}98 = 0{,}997248\), soit \(-0{,}2752\%\).
    \item \textbf{Défi C :} \(1{,}30 \times 0{,}70 = 0{,}91 \neq 1\). Le prix final est inférieur de \(9\%\) au prix initial.
\end{itemize}

\newpage

% ============================================================
\section*{Manche 3 -- Duel de capitaines (10 min, 8 pts)}

\textit{Les 4 capitaines s’affrontent. Chaque bonne réponse rapporte 2 points à son équipe. Les autres membres peuvent aider mais seul le capitaine répond.}

\vspace{0.3cm}

\begin{carte}{orange}{Duel 1}
Un capital augmente de \(10\%\) la première année, puis de \(20\%\) la deuxième.\\[0.2cm]
\textbf{Quel taux unique aurait donné le même résultat en 2 ans ?}
\end{carte}

\begin{carte}{orange}{Duel 2}
Un prix a augmenté de \(44\%\) en deux ans, avec le même taux annuel \(t\).\\[0.2cm]
\textbf{Déterminer \(t\).}
\end{carte}

\begin{carte}{orange}{Duel 3}
Une grandeur a subi 3 évolutions successives : \(+10\%\), \(-20\%\), \(+25\%\).\\[0.2cm]
\textbf{Quel est le taux global ?}
\end{carte}

\begin{carte}{orange}{Duel 4}
Un prix passe de \(80\)~\euro{} à \(100\)~\euro{}.\\[0.2cm]
\textbf{Quel est le taux d’évolution ? Puis, quel taux faudrait-il appliquer pour revenir à \(80\)~\euro{} ?}
\end{carte}

\vspace{0.3cm}
\textbf{Corrigé manche 3 :}
\begin{itemize}
    \item Duel 1 : \(1{,}10 \times 1{,}20 = 1{,}32\), soit \(+32\%\).
    \item Duel 2 : \(\sqrt{1{,}44} = 1{,}2\), soit \(t = 20\%\).
    \item Duel 3 : \(1{,}10 \times 0{,}80 \times 1{,}25 = 1{,}10\), soit \(+10\%\).
    \item Duel 4 : \(t_1 = +25\%\). Retour : \(t_2 = \frac{80-100}{100} \times 100 = -20\%\).
\end{itemize}

\newpage

% ============================================================
\section*{Manche 4 -- Enveloppe mystère (10 min, 12 pts)}

\textit{Chaque équipe tire une enveloppe au hasard. Elle contient un problème « mystère ». Chaque enveloppe vaut 3 points.}

\vspace{0.3cm}

\begin{multicols}{2}

\begin{carte}{violet}{Enveloppe A}
Un commerçant augmente ses prix de \(20\%\). Pour compenser une baisse des ventes, il les diminue ensuite de \(20\%\).\\[0.2cm]
\textbf{Les prix sont-ils revenus à leur niveau initial ? Justifier.}
\end{carte}

\begin{carte}{violet}{Enveloppe B}
Un article coûte \(250\)~\euro{} après une augmentation de \(25\%\).\\[0.2cm]
\textbf{Quel était son prix avant l’augmentation ?}
\end{carte}

\begin{carte}{violet}{Enveloppe C}
Un placement a rapporté \(21\%\) en deux ans, avec le même taux annuel \(t\).\\[0.2cm]
\textbf{Déterminer \(t\) (arrondir à \(0{,}01\%\)).}
\end{carte}

\begin{carte}{violet}{Enveloppe D}
Une population augmente de \(5\%\) par an pendant 3 ans.\\[0.2cm]
\textbf{Quel est le taux global ? Quel est le taux annuel moyen ?}
\end{carte}

\end{multicols}

\vspace{0.3cm}
\textbf{Corrigé manche 4 :}
\begin{itemize}
    \item A : \(1{,}20 \times 0{,}80 = 0{,}96 \neq 1\). Baisse de \(4\%\).
    \item B : \(V_i = \frac{250}{1{,}25} = 200\)~\euro{}.
    \item C : \(\sqrt{1{,}21} = 1{,}10\), soit \(t = 10\%\).
    \item D : CM \(= 1{,}05^3 \approx 1{,}157625\), soit \(+15{,}76\%\). Taux moyen \(= 5\%\) (car \(1{,}05^3\) correspond à 3 hausses de \(5\%\)).
\end{itemize}

\newpage

% ============================================================
\section*{Manche 5 -- Finale « Tout ou rien » (10 min, 10 pts)}

\textit{Chaque équipe choisit un niveau de difficulté. Une bonne réponse rapporte les points indiqués ; une mauvaise réponse les retire. Les points peuvent devenir négatifs.}

\vspace{0.3cm}

\begin{center}
\renewcommand{\arraystretch}{1.5}
\begin{tabular}{|c|p{9cm}|c|c|}
\hline
\textbf{Niveau} & \textbf{Question} & \textbf{Gain} & \textbf{Perte} \\
\hline
Facile & Un prix augmente de \(30\%\). Quel coefficient appliquer pour revenir au prix initial ? (arrondir à \(0{,}001\)) & \(+2\) & \(-2\) \\
\hline
Moyen & Un capital de \(5\,000\)~\euro{} subit 4 évolutions : \(+8\%\), \(-5\%\), \(+12\%\), \(-10\%\). Quel est le capital final ? & \(+3\) & \(-3\) \\
\hline
Difficile & Une grandeur a subi 5 évolutions annuelles identiques de taux \(t\). Le taux global est \(+27{,}628\%\). Déterminer \(t\). & \(+5\) & \(-5\) \\
\hline
\end{tabular}
\end{center}

\vspace{0.3cm}
\textbf{Corrigé manche 5 :}
\begin{itemize}
    \item Facile : \(CM = 1{,}30\), retour \(= \frac{1}{1{,}30} \approx 0{,}769\), soit \(-23{,}1\%\).
    \item Moyen : \(5\,000 \times 1{,}08 \times 0{,}95 \times 1{,}12 \times 0{,}90 \approx 5\,170{,}66\)~\euro{}.
    \item Difficile : \((1+t)^5 = 1{,}27628\), donc \(1+t = 1{,}05\), soit \(t = 5\%\).
\end{itemize}

\newpage

% ============================================================
\section*{Fiche de score}

\begin{center}
\renewcommand{\arraystretch}{1.8}
\begin{tabular}{|c|c|c|c|c|}
\hline
\textbf{Manche} & \textbf{Équipe 1} & \textbf{Équipe 2} & \textbf{Équipe 3} & \textbf{Équipe 4} \\
\hline
1 -- Échauffement (8 pts) & & & & \\
\hline
2 -- Défis (12 pts) & & & & \\
\hline
3 -- Duel (8 pts) & & & & \\
\hline
4 -- Enveloppes (12 pts) & & & & \\
\hline
5 -- Finale (10 pts) & & & & \\
\hline
\textbf{Total (50 pts)} & & & & \\
\hline
\end{tabular}
\end{center}

\vspace{0.5cm}

\begin{regle}
\textbf{Conseils pour l’enseignant :}
\begin{itemize}
    \item Distribuer la fiche de score dès le début.
    \item Chronométrer chaque manche rigoureusement pour respecter les 55 minutes.
    \item Pour la manche 1, utiliser un buzzer ou un simple lever de main.
    \item Pour la manche 5, annoncer les points avant la réponse pour créer du suspense.
    \item En cas d’égalité, poser une question bonus : « Un prix augmente de \(100\%\), puis baisse de \(50\%\). Revient-il au prix initial ? » (Réponse : oui, car \(2 \times 0{,}5 = 1\).)
\end{itemize}
\end{regle}

\vspace{0.5cm}

\begin{center}
{\Large\bfseries\color{bleufonce} Bon jeu à toutes les équipes !}
\end{center}

\end{document}